\begin{table}
    \setlength{\abovecaptionskip}{0.1cm}
    \setlength{\belowcaptionskip}{-0.5cm}
    \centering
    \small
    \begin{tabular}{l|c|c} 
        \toprule[1.1pt]
        \textbf{Model} & \textbf{Script Acc.} & \textbf{Blind Acc.}\\
        \midrule[1.1pt]
        GPT-4o        & 0.04 & 0.62\\
        GPT-3.5-turbo & 0.07 & 0.51\\
        \hline
        Mistral-7b & 0.05 & 0.56\\
        Llama-2-7b & 0.06 & 0.35\\
        Qwen2.5-7b & 0.03 & 0.55\\
        Llama-3-8b & 0.04 & 0.54\\
        Llama-2-13b & 0.04 & 0.35\\
        Qwen2.5-14b & 0.06 & 0.61\\
        Llama-2-70b & 0.06 & 0.40\\
        Llama-3-70b & 0.04 & 0.59\\
        Qwen2.5-72b & 0.04 & 0.60\\
        \bottomrule[1.1pt]
    \end{tabular}
    \caption{Model performance on script prediction and blind test. The low \textbf{Script Acc.} indicates the model barely discerns the scripts, and the \textbf{Blind Acc.} establishes a baseline for model's reasoning ability.}
    \label{tab:data_leakage}
\end{table}

Since our data originates from scripts, it is essential to prevent data leakage and evaluate the effectiveness of leakage mitigation.
Data leakage may arise from 1) the model's prior knowledge of the scenario's script, and 2) the information provided by the scenario that could help infer others' private information.
We quantify this risk through two experiments:
(1) \textbf{Script Prediction}: Whether models can guess the original scripts from background information, with 245 test samples (one per template).
(2) \textbf{Blind Test}: Whether models can answer private information reasoning questions with initial scenario information before interactions, with 100 test questions asked three times.

The script prediction results in Table~\ref{tab:data_leakage} indicate that models are nearly unable to infer the original script from the background information, proving the effectiveness of leakage mitigation.
The blind test results also establish a baseline for each model's private information reasoning ability.
